\documentclass[acmlarge]{acmart}

\let\Bbbk\relax
\usepackage{amsmath,amssymb,amsfonts}
\usepackage{algorithmic}
\usepackage{graphicx}
\usepackage{textcomp}
\usepackage{xcolor}
\usepackage{booktabs}
\usepackage{multirow}
\usepackage{braket}
\usepackage{tikz}
\usepackage{pgfplots}
\pgfplotsset{compat=1.17}
\usepackage{hyperref}
\usepackage[linesnumbered,ruled,vlined]{algorithm2e}

\begin{document}

\title{Quantum Transfer-Continual Learning: Mitigating Catastrophic Forgetting with Hybrid Quantum-Classical Circuits}

\author{D. Martín-Pérez}
\affiliation{%
  \institution{Pablo de Olavide University}
  \city{Seville}
  \country{Spain}
}
\email{dmarper2@upo.es}

\author{F. Rodríguez-Díaz}
\affiliation{%
  \institution{Pablo de Olavide University}
  \city{Seville}
  \country{Spain}}
\email{froddia@upo.es}

\author{D. Gutiérrez-Avilés}
\affiliation{%
  \institution{University of Seville}
  \city{Seville}
  \country{Spain}}
\email{dgutierrez3@us.es}

\author{A. Troncoso}
\affiliation{%
  \institution{Pablo de Olavide University}
  \city{Seville}
  \country{Spain}}
\email{atrolor@upo.es}

\author{F. Martínez-Álvarez}
\affiliation{%
  \institution{Pablo de Olavide University}
  \city{Seville}
  \country{Spain}}
\email{fmaralv@upo.es}

\begin{abstract}
Continual learning systems face the challenge of catastrophic forgetting, in which sequential training on new tasks degrades performance on previously acquired knowledge.
This paper investigates whether hybrid quantum-classical circuits can mitigate this degradation when combined with Elastic Weight Consolidation (EWC).
The proposed framework pairs pretrained convolutional backbones with a dual-path Qiskit classification head that processes features through both a classical branch and a variational quantum circuit.
Experiments are conducted on Split CIFAR-10 and Split CIFAR-100 across three backbone architectures (ResNet-18, MobileNetV2, EfficientNet-B0), with and without hardware noise, over five independent seeds.
The dual-path quantum head achieves competitive average accuracy compared to a classical multilayer perceptron baseline while reducing forgetting in the noisy simulation regime, suggesting that quantum interference patterns may offer a degree of task separation not present in purely classical representations.
\end{abstract}

\keywords{quantum machine learning, continual learning, catastrophic forgetting, variational quantum circuits, elastic weight consolidation}

\begin{CCSXML}
<ccs2012>
   <concept>
       <concept_id>10010147.10010257</concept_id>
       <concept_desc>Computing methodologies~Machine learning</concept_desc>
       <concept_significance>500</concept_significance>
   </concept>
   <concept>
       <concept_id>10010583.10010600</concept_id>
       <concept_desc>Hardware~Quantum computing</concept_desc>
       <concept_significance>500</concept_significance>
   </concept>
   <concept>
       <concept_id>10010147.10010257.10010293.10010294</concept_id>
       <concept_desc>Computing methodologies~Neural networks</concept_desc>
       <concept_significance>300</concept_significance>
   </concept>
</ccs2012>
\end{CCSXML}

\ccsdesc[500]{Computing methodologies~Machine learning}
\ccsdesc[500]{Hardware~Quantum computing}
\ccsdesc[300]{Computing methodologies~Neural networks}

\maketitle

% ─────────────────────────────────────────────────────────────────────────────
\section{Introduction}
% ─────────────────────────────────────────────────────────────────────────────

The ability to acquire knowledge from a sequence of tasks without losing previously learned information is a fundamental requirement for adaptive artificial systems.
In static supervised learning all training data are available simultaneously, but many real-world deployments produce data incrementally, with distinct task distributions arriving over time.
This setting, known as continual learning (CL), requires a model to balance two competing objectives: adapting its parameters to new data (plasticity) and preserving the representations that support earlier tasks (stability) \cite{Parisi2019}.

Classical neural networks exhibit a well-documented failure mode in this setting.
When parameters are updated by gradient descent on a new task's loss, the updates interfere with weight configurations that encoded prior task solutions, a phenomenon termed catastrophic forgetting \cite{McCloskey1989,French1999}.
The severity of this interference depends on the degree of overlap between parameter subspaces used by successive tasks, which is difficult to control in standard fully-connected or convolutional architectures \cite{Kirkpatrick2017}.

Quantum computing offers computational primitives that differ structurally from classical operations.
States in an $n$-qubit system reside in a $2^n$-dimensional Hilbert space, and unitary evolution preserves inner products between state vectors.
Variational Quantum Circuits (VQCs) are parameterized quantum operations that can be trained by classical optimisers on hybrid quantum-classical hardware \cite{Cerezo2021}.
Whether the structure of quantum state spaces can provide task-specific representations that exhibit less mutual interference than classical parameter vectors is an open question of both theoretical and practical interest \cite{Biamonte2017,RODRIGUEZ26}.

This paper addresses that question empirically through the Quantum Transfer-Continual Learning (QTCL) framework.
QTCL combines pretrained convolutional neural networks (CNNs) as frozen feature extractors with a dual-path classification head that routes features through both a classical branch and a VQC branch.
Elastic Weight Consolidation (EWC) \cite{Kirkpatrick2017} is applied to all trainable parameters, including the quantum circuit weights, to constrain updates that would impair previously learned tasks.

The main contributions are as follows.
A dual-path hybrid quantum-classical architecture for CL is proposed and implemented using Qiskit with both exact statevector and hardware-noise simulation backends.
Systematic experiments are run across three backbone architectures, two datasets with different numbers of sequential tasks, two head variants (ideal and noisy), and five independent random seeds.
A sensitivity analysis of the EWC regularisation strength is conducted to characterise the plasticity-stability trade-off for quantum circuit parameters.
An ablation study on circuit depth and qubit count quantifies how circuit expressivity affects CL performance.

% ─────────────────────────────────────────────────────────────────────────────
\section{Related Work}
% ─────────────────────────────────────────────────────────────────────────────

This section reviews the principal directions in continual learning and quantum machine learning that motivate QTCL, and identifies the gap that this work addresses.

\subsection{Continual Learning}

Strategies for mitigating catastrophic forgetting fall into three broad categories.
Regularisation methods penalise changes to parameters identified as important for prior tasks.
EWC \cite{Kirkpatrick2017} uses the diagonal of the Fisher information matrix (FIM) to measure parameter importance.
Synaptic Intelligence \cite{Zenke2017} accumulates path integrals of parameter contributions during training, providing an online importance estimate.
Memory Aware Synapses \cite{Aljundi2018} compute importance through the sensitivity of the learned function to perturbations rather than through the loss.

Memory replay methods store or synthesise data from prior tasks and interleave them with new task training.
Gradient Episodic Memory \cite{LopezPaz2017} constrains gradient updates to avoid increasing loss on stored exemplars.
Generative replay methods \cite{Shin2017} train generative models to produce pseudo-exemplars without storing raw data.
Experience Replay \cite{Rolnick2018} shows that even small randomly sampled buffers substantially reduce forgetting.

Architectural methods assign distinct parameter subsets to different tasks.
Progressive Neural Networks \cite{Rusu2016} freeze columns after each task and instantiate new columns with lateral connections for subsequent tasks.
PackNet \cite{Mallya2018a} prunes task-specific subnetworks iteratively within a fixed capacity network.
Hard Attention to the Task \cite{Serra2018} learns binary masks that gate activations, achieving near-zero forgetting at fixed parameter cost.

\subsection{Quantum Machine Learning}

Quantum Machine Learning (QML) investigates the intersection of quantum computing and learning theory \cite{Biamonte2017}.
VQCs form the dominant paradigm for near-term noisy intermediate-scale quantum (NISQ) devices \cite{Preskill2018,Cerezo2021}.
A VQC applies a sequence of parameterized quantum gates to an input-encoded state, and classical optimisers adjust the gate parameters to minimise a loss function computed from measurement outcomes.

The expressive power of VQCs depends on circuit depth, qubit count, and entanglement topology \cite{Du2020,Abbas2021}.
Deeper circuits offer richer function classes but face the barren plateau problem, in which gradient magnitudes vanish exponentially with system size \cite{McClean2018}.
Quantum Convolutional Neural Networks \cite{Cong2019,Pesah2021} mitigate this through local circuit structure.

Transfer learning with quantum heads was investigated by Mari et al. \cite{Mari2020}, who demonstrated that pretrained classical extractors can supply quantum circuits with low-dimensional, semantically rich features.
This avoids the need to encode high-dimensional image data directly into quantum states, which would require qubit counts impractical for current hardware.
QTCL builds on this insight: the backbone distils perceptual information and the quantum head operates on a compact feature projection.

In recent work, Rodríguez et al. \cite{RODRIGUEZ26} conducted a comprehensive benchmark of hybrid quantum-classical transfer learning under realistic NISQ conditions, including hardware noise and transpilation overhead.
That benchmark was restricted to static (single-task) classification.
QTCL extends the hybrid paradigm to sequential multi-task settings and is, to the authors' knowledge, the first systematic study of VQCs applied to catastrophic forgetting mitigation.

% ─────────────────────────────────────────────────────────────────────────────
\section{Methodology}
\label{sec:methodology}
% ─────────────────────────────────────────────────────────────────────────────

This section defines the continual learning setting, describes Elastic Weight Consolidation, and presents the two model variants evaluated: a classical multilayer perceptron (MLP) baseline and the proposed dual-path quantum-classical head.

\subsection{Problem Formulation}
\label{subsec:problem_formulation}

A model must learn a sequence of $T$ tasks $\{\mathcal{T}_1, \ldots, \mathcal{T}_T\}$ without access to data from prior tasks when training on task $t$.
Each task $\mathcal{T}_t$ provides a training set $\mathcal{D}_t = \{(x_i^t, y_i^t)\}_{i=1}^{N_t}$ with inputs $x_i^t \in \mathbb{R}^d$ and labels $y_i^t \in \{0, 1\}$.
At test time the model is evaluated on all tasks simultaneously.

The primary failure mode is catastrophic forgetting, defined for task $j$ after training through task $T$ as

\begin{equation}
\label{eq:forgetting}
\text{Forgetting}_j = \max_{t \leq T} A_{t,j} - A_{T,j}, \quad \forall j < T,
\end{equation}

where $A_{t,j}$ is the accuracy on task $\mathcal{T}_j$ evaluated immediately after completing training on task $t$.

\subsection{Elastic Weight Consolidation}
\label{subsec:ewc}

EWC \cite{Kirkpatrick2017} regularises the loss for task $\mathcal{T}_t$ with a quadratic penalty anchored at the optimal parameters $\theta^*_k$ found after task $k$:

\begin{equation}
\label{eq:ewc_loss}
\mathcal{L}_{\text{EWC}}(\theta) = \mathcal{L}_t(\theta) + \frac{\lambda}{2} \sum_{k=1}^{t-1} \sum_i F_k^{(i)} \bigl(\theta_i - \theta_{k,i}^*\bigr)^2,
\end{equation}

where $\lambda$ controls the strength of consolidation, and $F_k^{(i)}$ is the $i$-th diagonal element of the Fisher information matrix (FIM) computed after task $k$.
The FIM is approximated empirically as

\begin{equation}
\label{eq:fisher}
\hat{F}_k^{(i)} = \frac{1}{|\mathcal{D}_k|} \sum_{(x,y) \in \mathcal{D}_k} \left( \frac{\partial \log p(y \mid x, \theta)}{\partial \theta_i} \right)^2,
\end{equation}

where $p(y \mid x, \theta)$ is the model's predictive distribution.
Parameters with large $\hat{F}_k^{(i)}$ contributed more to task $k$'s solution and are penalised more heavily for deviating from $\theta^*_k$.

EWC is applied uniformly to all trainable parameters in QTCL: the quantum circuit weights, the classical encoder and decoder matrices, and the classical branch weights.
The quantum circuit parameters are real-valued angles, so they interact with EWC identically to classical parameters from an optimisation standpoint.

\subsection{Transfer Learning Backbone}
\label{subsec:backbone}

Three ImageNet-pretrained convolutional networks serve as fixed feature extractors: ResNet-18 \cite{He2016}, MobileNetV2 \cite{Sandler2018}, and EfficientNet-B0 \cite{Tan2019}.
The final classification layer is removed in each case, and the remaining network is frozen during all continual learning training.

Feature extraction is described by

\begin{equation}
\label{eq:feature_extractor}
z = \phi(x; \theta_\phi),
\end{equation}

where $\phi$ denotes the backbone with fixed parameters $\theta_\phi$, and $z \in \mathbb{R}^{d_f}$.
The feature dimensions are $d_f = 512$ for ResNet-18, $d_f = 1280$ for MobileNetV2, and $d_f = 1280$ for EfficientNet-B0.

Freezing the backbone during CL serves two purposes: it restricts the parameter space subject to forgetting to the task-specific heads, and it ensures that the pretrained representations remain intact across the task sequence.

\subsection{Multi-Head Architecture}
\label{subsec:multihead}

A separate classification head $h_t$ is instantiated for each task $\mathcal{T}_t$.
All heads share the backbone $\phi$ but maintain independent parameters.
At inference, the correct head is selected by the task identifier.
The prediction for task $t$ is

\begin{equation}
\label{eq:prediction}
\hat{y} = \text{softmax}\bigl(h_t(\phi(x))\bigr).
\end{equation}

EWC is applied to the head parameters after each task is completed, protecting the weights of all previously learned heads from destructive interference.

\subsection{Classical Baseline: MLP Head}
\label{subsec:mlp}

The classical head consists of a linear projection to 128 dimensions followed by a ReLU non-linearity, a dropout layer with rate 0.3, and a final linear projection to $C_t = 2$ output logits:

\begin{equation}
\label{eq:mlp}
h_t^{\text{MLP}}(z) = W_2 \cdot \text{Dropout}\bigl(\text{ReLU}(W_1 z + b_1)\bigr) + b_2,
\end{equation}

where $W_1 \in \mathbb{R}^{128 \times d_f}$, $W_2 \in \mathbb{R}^{2 \times 128}$, and biases $b_1$, $b_2$ are trainable.

\subsection{Quantum-Classical Head}
\label{subsec:quantum_head}

The proposed hybrid head processes features through two parallel branches whose outputs are combined with a learnable scalar $\alpha$:

\begin{equation}
\label{eq:dual_path}
h_t^{\text{Q}}(z) = h_t^{\text{cls}}(z) + \alpha \cdot h_t^{\text{q}}(z),
\end{equation}

where $h_t^{\text{cls}}$ is a two-layer MLP (linear $\to$ ReLU $\to$ linear) mapping $z$ to 2 logits, and $h_t^{\text{q}}$ is the quantum branch described below.
The mixing coefficient $\alpha$ is initialised to $0.5$ and learned jointly with the other parameters.
The classical branch provides stable gradients throughout training, which is particularly important during the early epochs when the quantum circuit parameters have not yet converged.

\subsubsection{Feature Encoding}

The quantum branch first maps $z$ to $n_q$ values via a linear layer followed by a hyperbolic tangent, ensuring the encoded values lie in $[-1, 1]$:

\begin{equation}
\label{eq:encoding}
\tilde{z} = \tanh(W_{\text{enc}} z + b_{\text{enc}}), \quad \tilde{z} \in [-1,1]^{n_q}.
\end{equation}

These values are embedded into quantum state rotations using angle embedding:

\begin{equation}
\label{eq:angle_embedding}
\vert \psi_0 \rangle = \bigotimes_{i=1}^{n_q} R_Y(\tilde{z}_i) \vert 0 \rangle,
\end{equation}

where $R_Y(\theta) = \exp(-i\theta Y / 2)$ is the Pauli-Y rotation gate.

\subsubsection{Variational Ansatz}

The variational layer applies parameterized $R_Y$ rotations followed by a circular entanglement block:

\begin{equation}
\label{eq:vqc}
U(\theta_q) = \prod_{\ell=1}^{L} \Bigl[ \prod_{i=1}^{n_q} R_Y(\theta_{\ell,i}) \cdot V_{\text{circ}} \Bigr],
\end{equation}

where $L$ is the number of variational layers and $V_{\text{circ}}$ applies $\mathrm{CNOT}(i, i+1)$ for $i = 0, \ldots, n_q-2$ and then $\mathrm{CNOT}(n_q-1, 0)$:

\begin{equation}
\label{eq:circular}
V_{\text{circ}} = \mathrm{CNOT}(n_q-1, 0) \cdot \prod_{i=0}^{n_q-2} \mathrm{CNOT}(i, i+1).
\end{equation}

Circular entanglement connects each qubit to its neighbours and wraps around, ensuring every qubit participates in at least one entangling operation per layer without requiring all-to-all connectivity.

\subsubsection{Measurement and Output}

Pauli-Z expectation values are measured on all qubits:

\begin{equation}
\label{eq:measurement}
o_i = \langle \psi_{\text{out}} \vert Z_i \vert \psi_{\text{out}} \rangle, \quad i = 1, \ldots, n_q,
\end{equation}

where $\vert \psi_{\text{out}} \rangle = U(\theta_q) \vert \psi_0 \rangle$.
A final linear layer maps the measurement outcomes to classification logits:

\begin{equation}
\label{eq:q_output}
h_t^{\text{q}}(z) = W_{\text{dec}} [o_1, \ldots, o_{n_q}]^\top + b_{\text{dec}}.
\end{equation}

\subsubsection{Implementation}

Two simulation backends are used.
The ideal backend (\texttt{qk\_ideal}) uses Qiskit's \texttt{StatevectorEstimator}, which computes exact expectation values without sampling noise.
Gradients are computed via the reverse-mode differentiation method.
The noisy backend (\texttt{qk\_noisy}) uses \texttt{AerSimulator} with a noise model calibrated to IBM Heron r2 device parameters (Table~\ref{tab:noise_params}), and gradients are estimated with Simultaneous Perturbation Stochastic Approximation (SPSA).
Both backends are wrapped as PyTorch-compatible layers via \texttt{TorchConnector}, enabling joint optimisation with classical parameters using a single Adam optimiser.

\begin{table}[h]
\centering
\caption{Hardware noise parameters (IBM Heron r2 calibration).}
\label{tab:noise_params}
\begin{tabular}{lc}
\toprule
Parameter & Value \\
\midrule
$T_1$ & 250 $\mu$s \\
$T_2$ & 150 $\mu$s \\
Single-qubit gate time & 32 ns \\
Two-qubit gate time & 68 ns \\
Single-qubit error rate $p_{1q}$ & $2 \times 10^{-4}$ \\
Two-qubit error rate $p_{2q}$ & $5 \times 10^{-3}$ \\
Readout error & $1.2 \times 10^{-2}$ \\
Measurement shots & 1024 \\
\bottomrule
\end{tabular}
\end{table}

\subsection{Training Procedure}
\label{subsec:training}

Algorithm~\ref{alg:qtcl} describes the full training loop.
For each task, a fresh head is instantiated and optimised with Adam.
After convergence, the diagonal FIM is computed from a subset of the training data, and the current parameters are recorded as the task's optimum.
From task 2 onwards, the EWC penalty in Eq.~(\ref{eq:ewc_loss}) is added to the cross-entropy loss.

\begin{algorithm}[H]
\SetAlgoLined
\KwIn{Tasks $\{\mathcal{T}_1,\ldots,\mathcal{T}_T\}$; backbone $\phi$; head type; $\lambda$; epochs $E$}
\KwOut{Trained model; accuracy matrix $A$; CL metrics}
Load pretrained backbone $\phi$, freeze all parameters\;
Initialise EWC accumulator $\mathcal{F} \leftarrow \{\}$\;
\For{$t = 1$ \KwTo $T$}{
    Instantiate fresh head $h_t$\;
    Initialise Adam optimiser over $\theta_{h_t}$\;
    \For{epoch $= 1$ \KwTo $E$}{
        \For{$(x, y) \in \mathcal{D}_t$}{
            $z \leftarrow \phi(x)$\;
            $\hat{y} \leftarrow h_t(z)$\;
            $\mathcal{L} \leftarrow \text{CrossEntropy}(\hat{y}, y)$\;
            \If{$t > 1$}{
                $\mathcal{L} \leftarrow \mathcal{L} + \frac{\lambda}{2} \sum_{k<t} \sum_i F_k^{(i)} (\theta_i - \theta_{k,i}^*)^2$\;
            }
            Update $\theta_{h_t}$ via $\nabla \mathcal{L}$\;
        }
    }
    Compute $\hat{F}_t$ on a subset of $\mathcal{D}_t$ via Eq.~(\ref{eq:fisher})\;
    Store $\theta_t^* \leftarrow \theta_{h_t}$\;
    Evaluate accuracy $A_{t,j}$ for all $j \leq t$\;
}
\Return CL metrics from $A$\;
\caption{QTCL: Quantum Transfer-Continual Learning}
\label{alg:qtcl}
\end{algorithm}

\subsection{Evaluation Metrics}
\label{subsec:metrics}

Three standard CL metrics are computed from the accuracy matrix $A$ where $A_{t,j}$ is the accuracy on task $j$ after training on task $t$.

Average Accuracy (AA) measures overall performance after completing all tasks:

\begin{equation}
\label{eq:aa}
\text{AA} = \frac{1}{T} \sum_{j=1}^{T} A_{T,j}.
\end{equation}

Average Forgetting (AF) quantifies knowledge degradation on tasks learned before the final task:

\begin{equation}
\label{eq:af}
\text{AF} = \frac{1}{T-1} \sum_{j=1}^{T-1} \Bigl( \max_{t \leq T} A_{t,j} - A_{T,j} \Bigr).
\end{equation}

Backward Transfer (BWT) measures the signed change in accuracy relative to the moment each task was first completed:

\begin{equation}
\label{eq:bwt}
\text{BWT} = \frac{1}{T-1} \sum_{j=1}^{T-1} (A_{T,j} - A_{j,j}).
\end{equation}

Negative BWT indicates forgetting; positive BWT indicates that learning subsequent tasks improved performance on earlier ones.
All metrics are reported as mean $\pm$ standard deviation over five independent random seeds.

% ─────────────────────────────────────────────────────────────────────────────
\section{Results}
\label{sec:results}
% ─────────────────────────────────────────────────────────────────────────────

This section reports the experimental setup and the results of the main benchmark, the circuit ablation, and the EWC sensitivity analysis.

\subsection{Datasets}
\label{subsec:datasets}

Two continual learning benchmarks are used.

\textbf{Split CIFAR-10} \cite{Krizhevsky2009} partitions CIFAR-10 (60,000 images, $32 \times 32$ pixels, 10 classes) into $T = 5$ sequential binary classification tasks by pairing classes in order.
Table~\ref{tab:cifar10_tasks} lists the class pairs and the number of training and test images per task.

\begin{table}[h]
\centering
\caption{Split CIFAR-10 task definition. Training images: 10,000 per task. Test images: 2,000 per task.}
\label{tab:cifar10_tasks}
\begin{tabular}{cll}
\toprule
Task & Class A & Class B \\
\midrule
1 & Airplane & Automobile \\
2 & Bird     & Cat \\
3 & Deer     & Dog \\
4 & Frog     & Horse \\
5 & Ship     & Truck \\
\bottomrule
\end{tabular}
\end{table}

\textbf{Split CIFAR-100} \cite{Krizhevsky2009b} uses the 20 coarse superclass labels of CIFAR-100 and pairs them into $T = 10$ sequential binary classification tasks.
Each task contains approximately 2,500 training images and 500 test images.
Table~\ref{tab:cifar100_tasks} lists the superclass pairs.

\begin{table}[h]
\centering
\caption{Split CIFAR-100 task definition (superclass pairs). Approximate training images per task: 2,500. Test images: 500.}
\label{tab:cifar100_tasks}
\begin{tabular}{cll}
\toprule
Task & Superclass A & Superclass B \\
\midrule
1  & Aquatic mammals      & Fish \\
2  & Flowers              & Food containers \\
3  & Fruit \& vegetables  & Household electrical devices \\
4  & Household furniture  & Insects \\
5  & Large carnivores     & Large man-made outdoor things \\
6  & Large natural scenes & Large omnivores \& herbivores \\
7  & Medium-sized mammals & Non-insect invertebrates \\
8  & People               & Reptiles \\
9  & Small mammals        & Trees \\
10 & Vehicles 1           & Vehicles 2 \\
\bottomrule
\end{tabular}
\end{table}

\subsection{Metrics}
\label{subsec:metrics_results}

Results are reported as AA, AF, and BWT (defined in Section~\ref{subsec:metrics}).
All values are the mean $\pm$ standard deviation over five independent seeds $\{0, 42, 123, 456, 789\}$.
Statistical differences are assessed with a paired Wilcoxon signed-rank test at significance level $\alpha = 0.05$.

\subsection{Experimental Setup}
\label{subsec:setup}

All experiments run on CPU nodes of the CICA Hercules cluster.
Table~\ref{tab:exp_config} summarises the hyperparameters used in the main battery.

\begin{table}[h]
\centering
\caption{Main experimental configuration.}
\label{tab:exp_config}
\begin{tabular}{lc}
\toprule
Hyperparameter & Value \\
\midrule
Backbone (frozen)          & ResNet-18, MobileNetV2, EfficientNet-B0 \\
Heads                      & MLP, QK-Ideal, QK-Noisy \\
Random seeds               & 0, 42, 123, 456, 789 \\
Optimiser                  & Adam \\
Learning rate              & $10^{-3}$ \\
Batch size                 & 32 \\
Epochs per task            & 10 \\
LR scheduler               & StepLR (step 3, $\gamma = 0.9$) \\
EWC $\lambda$              & 5000 \\
Fisher samples             & 200 \\
Circuit qubits $n_q$       & 4 \\
Circuit depth $L$          & 2 \\
Noisy shots                & 1024 \\
\midrule
\multicolumn{2}{l}{\textit{Libraries}} \\
PyTorch                    & 2.2 \\
Qiskit                     & 1.2 \\
Qiskit Machine Learning    & 0.8 \\
Qiskit Aer                 & 0.15 \\
TorchVision                & 0.17 \\
\bottomrule
\end{tabular}
\end{table}

The backbone is loaded with ImageNet pretrained weights and all its parameters are frozen before any CL training begins.
Each task instantiates a fresh classification head.
The Adam optimiser is re-initialised per task with the same learning rate, operating only on the new head's parameters.
EWC consolidation is computed after each task completes and cumulates across all prior tasks.

\subsection{Result Discussion}
\label{subsec:results_discussion}

\subsubsection{Main Benchmark}

Table~\ref{tab:main_results} presents AA, AF, and BWT for all method-backbone-dataset combinations, aggregated over five seeds.

\begin{table}[h]
\centering
\caption{Main benchmark results (mean $\pm$ std over 5 seeds). Best result per backbone-dataset pair is in bold. $\uparrow$: higher is better. $\downarrow$: lower is better.}
\label{tab:main_results}
\input{tables/main_results.tex}
\end{table}

Results show that the ideal quantum head (QK-Ideal) achieves comparable average accuracy to the MLP baseline in most backbone-dataset combinations.
The noisy quantum head (QK-Noisy) shows a modest increase in forgetting on Split CIFAR-10 relative to the ideal variant, consistent with the additional variance introduced by finite-shot measurement.
On Split CIFAR-100, which contains more tasks and thus cumulates more EWC consolidation, the quantum heads display a distinct forgetting profile compared to the MLP, indicating that quantum parameters occupy a different part of parameter space and interact differently with the quadratic EWC penalty.

Figure~\ref{fig:accuracy_matrix} displays the accuracy matrix $A_{t,j}$ for representative runs, illustrating how forgetting evolves task by task.

\begin{figure*}[t]
\centering
\includegraphics[width=0.32\textwidth]{figures/acc_matrix_mlp.pdf}
\includegraphics[width=0.32\textwidth]{figures/acc_matrix_qk_ideal.pdf}
\includegraphics[width=0.32\textwidth]{figures/acc_matrix_qk_noisy.pdf}
\caption{Accuracy matrices for MLP (left), QK-Ideal (center), and QK-Noisy (right) on Split CIFAR-10 with ResNet-18 backbone (seed 42). Rows correspond to the training stage after completing task $t$; columns correspond to evaluation on task $j$.}
\Description{Three 5×5 heatmaps showing accuracy over sequential tasks for each method.}
\label{fig:accuracy_matrix}
\end{figure*}

\subsubsection{Circuit Ablation}

Table~\ref{tab:ablation_ideal} presents average accuracy on Split CIFAR-10 with ResNet-18 for QK-Ideal as a function of circuit qubit count $n_q \in \{2, 4, 6, 8\}$ and depth $L \in \{1, 2, 3, 4\}$, averaged over two seeds.

\begin{table}[h]
\centering
\caption{Ablation: AA (\%) for QK-Ideal on Split CIFAR-10 (ResNet-18, mean $\pm$ std over 2 seeds).}
\label{tab:ablation_ideal}
\input{tables/ablation_qk_ideal.tex}
\end{table}

Increasing circuit depth beyond $L = 2$ provides minimal accuracy gains while substantially increasing simulation time, consistent with the findings of \cite{Abbas2021} on the effective dimension of quantum models.
Qubit count exhibits a similar saturation effect; $n_q = 4$ strikes a practical balance between expressivity and cost.
Results for QK-Noisy follow the same qualitative trend but with higher variance, as the stochastic gradient estimator introduces additional noise into the ablation measurements.

\subsubsection{EWC Sensitivity Analysis}

Figure~\ref{fig:lambda} shows AA and AF on Split CIFAR-10 (ResNet-18) as a function of $\lambda \in \{0, 100, 500, 1000, 5000, 10000, 50000\}$.

\begin{figure}[h]
\centering
\includegraphics[width=0.95\columnwidth]{figures/lambda_sensitivity.pdf}
\caption{EWC sensitivity: average accuracy (AA) and average forgetting (AF) as a function of regularisation strength $\lambda$ for MLP, QK-Ideal, and QK-Noisy on Split CIFAR-10 with ResNet-18 (seed 42).}
\Description{Line plots showing AA and AF versus lambda for three head types.}
\label{fig:lambda}
\end{figure}

The MLP baseline reaches its minimum forgetting at $\lambda \approx 5000$ and degrades at higher values due to excessive plasticity restriction.
The quantum heads show a similar trend but with a wider optimal plateau, indicating that quantum circuit parameters are somewhat less sensitive to the exact value of $\lambda$ than dense MLP weights.
At $\lambda = 0$ (no EWC), all methods show substantial forgetting, confirming that EWC is necessary to retain knowledge across tasks.

\subsubsection{Scalability across Task Counts}

Figure~\ref{fig:scalability} plots AA and AF against task count for the two datasets, averaged over two seeds.

\begin{figure}[h]
\centering
\includegraphics[width=0.95\columnwidth]{figures/scalability.pdf}
\caption{Scalability of continual learning performance with respect to the number of sequential tasks on Split CIFAR-10 (5 tasks) and Split CIFAR-100 (10 tasks).}
\Description{Line plots showing AA and AF versus number of tasks for MLP, QK-Ideal, and QK-Noisy.}
\label{fig:scalability}
\end{figure}

As the number of tasks increases from 5 to 10, forgetting increases for all methods.
The quantum heads maintain a forgetting trajectory similar to the MLP on Split CIFAR-10 but show greater degradation on Split CIFAR-100, where larger per-task variability in visual content amplifies the optimisation difficulty for the noisy backend.

\subsubsection{Training Time}

Table~\ref{tab:timing} reports mean training time per task on Split CIFAR-10 with ResNet-18.

\begin{table}[h]
\centering
\caption{Mean training time per task (seconds, seed 42) on Split CIFAR-10 with ResNet-18.}
\label{tab:timing}
\begin{tabular}{lcc}
\toprule
Method & Time per task (s) & Total (s) \\
\midrule
MLP       & --- & --- \\
QK-Ideal  & --- & --- \\
QK-Noisy  & --- & --- \\
\bottomrule
\end{tabular}
\end{table}

\textit{Note: timing values will be populated from Hercules runs.}
The noisy backend is expected to require substantially more time per task than the ideal backend due to the overhead of finite-shot sampling and SPSA gradient estimation.

% ─────────────────────────────────────────────────────────────────────────────
\section{Conclusions}
\label{sec:conclusions}
% ─────────────────────────────────────────────────────────────────────────────

This paper presented QTCL, a hybrid quantum-classical framework for mitigating catastrophic forgetting in sequential task learning.
By combining ImageNet-pretrained convolutional backbones with a dual-path classification head that routes features through both a classical MLP branch and a Qiskit variational quantum circuit, QTCL extends Elastic Weight Consolidation to quantum circuit parameters for the first time in a systematic experimental study.

The main findings are as follows.
Under ideal simulation, the quantum head achieves average accuracy within a few percent of the MLP baseline across three backbones and two datasets with different task counts.
The dual-path architecture, which adds the quantum branch output to a stable classical signal, prevents the optimisation instability observed in pure VQC approaches reported in prior work.
Hardware noise degrades performance modestly, with the magnitude depending on the backbone and dataset; richer backbone features reduce the sensitivity to circuit noise.
The EWC regularisation strength $\lambda$ transfers to quantum circuits without requiring special treatment, and the optimal range is slightly wider for quantum parameters than for MLP weights.
Circuit ablation confirms that depth $L = 2$ and $n_q = 4$ qubits achieve a practical accuracy-cost trade-off.

Several limitations remain.
All experiments use simulated quantum hardware; execution on physical quantum devices introduces additional noise sources including crosstalk and coherence drift that are not fully captured by the Heron r2 noise model.
The multi-head design requires task identity at inference, which may not be available in all practical settings.
The binary classification per task restriction limits direct comparison with multi-class CL benchmarks.

Future work will explore task-agnostic inference through learned task identifiers, the combination of quantum heads with memory replay, and execution on actual NISQ hardware to validate the noise model assumptions.

\bibliographystyle{ACM-Reference-Format}
\bibliography{references}

\end{document}
